\documentclass[10pt]{letter}
\usepackage[T1]{fontenc}
\usepackage[utf8]{inputenc}
\usepackage{lmodern}
\usepackage{microtype}
\usepackage[margin=0.72in]{geometry}
\usepackage{hyperref}
\hypersetup{hidelinks}
\signature{Thaisse Dias Paes\\Corresponding author\\Federal University of Par\'a\\\href{mailto:thaisse.paes@itec.ufpa.br}{thaisse.paes@itec.ufpa.br}}
\address{Electrical Engineering Graduate Program\\Federal University of Par\'a\\Bel\'em, Par\'a, Brazil}
\date{30 August 2026}

\begin{document}
\begin{letter}{Editors\\Journal of Nonlinear Science}
\opening{Dear Editors,}

Please consider our manuscript, ``Affine-invariant isochron curvature predicts finite-amplitude phase-reduction error in nonlinear oscillators,'' for publication in the \textit{Journal of Nonlinear Science}.

The paper addresses a general problem in nonlinear model reduction: first-order phase reduction is local, yet the curvature controlling its leading finite-amplitude error is coordinate dependent. We show that covariance whitening converts every nonsingular affine reparameterization of the state into an orthogonal change of whitened coordinates, making norm- and spectrum-based local curvature descriptors invariant. We then test the fitted phase gradient and Hessian prospectively across 63 oscillator settings using perturbation directions disjoint from those used for derivative estimation.

The numerical results closely follow the theoretical prediction. The first-order residual scales with exponent 2.0017, the Hessian-corrected residual with exponent 3.0255, and at normalized amplitude $\varepsilon=0.10$ the fitted quadratic term predicts the signed first-order residual with median slope 0.9980 and median $R^2=0.9834$. The Hessian reduces the setting-level median phase-reset error by 94.1\% at $\varepsilon=0.10$ and 83.8\% at $\varepsilon=0.25$, with positive improvement in all 63 settings. A 15,750-transformation affine stress test shows large coordinate dependence in native and independent unit-box descriptors while covariance-whitened descriptors remain equivalent.

A shorter conference paper accepted for presentation at CBEB 2026 examined descriptive isochron-foliation fingerprints and model-family separation. The present manuscript addresses a distinct question and develops the affine-equivalence result, time-phase covariance whitening, held-out finite-amplitude prediction, exact Hessian validation, and setting-level error analysis required for that question. The manuscript is not under consideration elsewhere, and all authors have approved its submission.

The submission includes Supplemental Material, setting-level numerical data, and figure-reproduction code. In accordance with the journal's request for suggested editors, we suggest Jeff Moehlis, Kyle Wedgwood, or Christian Bick because their expertise is closely aligned with nonlinear oscillators, phase reduction, and dynamical-systems methods. The work fits the \textit{Journal of Nonlinear Science} because its primary contribution is a general geometric and predictive result for nonlinear oscillator reduction.

\closing{Sincerely,}
\end{letter}
\end{document}
